%! TeX root = ../report.tex

\section{Introduction}
This paper is my results and methods of simulations regarding
pendulums, the language of C++ will be dear to us as we
proceed through each section. I will first give a review of
materials that we must know in order to model such
phenomenas. More specifically, we will start
with Hooke's law, which will introduce us to this notion
of Simple Harmonic Motion. I wrote this paper originally
as a way to gather my thoughts and build materials that 
I needed to model a double pendulum, which means this is
more informal that it is formal, but I will try my best to
use certain language to promote generality of terms, I also
wanted to make it reproduceable, so you, the reader, can
follow along.
\subsection{Prerequisites}
We should get this out of the way, you are going to need certain
tools on your device (performed on Windows, but can be performed
on other operating systems as long as you have the same tools)
\begin{center}
\begin{minipage}{0.78\textwidth}
\begin{enumerate}
    \item A modern C++ compiler supporting at least the C++17 standard (such as \texttt{g++}, \texttt{clang++}, or MSVC).
    \item A terminal or command-line interface to compile and execute the source files.
\end{enumerate}
\end{minipage}
\end{center}


